\subsection{RQ3: Paired-Consistency Anchoring}
\label{subsec:rq3}

\textbf{Experimental Setup.} We implement our paired-consistency anchoring objective, anchoring the student model to the clean-code teacher via KL divergence:
\begin{equation}
\mathcal{L} = \mathrm{CE}\!\left(y \mid x_{\mathrm{obf}}\right) + \lambda\, D_{\mathrm{KL}}\!\left( p_{\mathrm{teacher}}(\cdot \mid x_{\mathrm{clean}}) \,\Vert\, p_{\mathrm{student}}(\cdot \mid x_{\mathrm{obf}}) \right)
\end{equation}
We evaluate this against the Breadth Fine-tuning baseline from RQ1, focusing specifically on performance recovery on unseen obfuscation families and the retention of clean-code capabilities.

\textbf{Results and Analysis.} The KL-anchored model successfully preserves the gains of broad training while eliminating the penalty on unseen transformation families. As detailed in Table~\ref{tab:rq3_kl_results}, the anchored student achieves a substantial percentage point increase on unseen families compared to standard breadth training, without sacrificing accuracy on clean inputs.

\begin{table}[h]
\centering
\caption{Anchored Student vs. Baselines on Unseen Families}
\label{tab:rq3_kl_results}
\begin{tabular}{@{}lccc@{}}
\toprule
\textbf{Method}             & \textbf{Clean Code} & \textbf{Seen Stacks} & \textbf{Unseen Family} \\ \midrule
Teacher (Clean-tuned)       & 85.0\%              & 25.4\%               & 22.1\%                 \\
Breadth Fine-tuning         & 79.8\%              & 65.2\%               & 24.8\%                 \\
\textbf{Paired-Consistency} & \textbf{84.5\%}     & \textbf{66.1\%}      & \textbf{41.3\%}        \\ \bottomrule
\end{tabular}
\end{table}

By changing the learning target to the clean-code behavioral distribution, the model successfully abstracts away surface heuristics and relies on structural semantics.